\section{Kernel sums and the bounded guide fields}
\label{sec:seed-guides}
The energy argument uses kernels after parallel transport, truncation
and periodization. We derive their row and column bounds from the
cover estimates, and then bound the clipped guide fields. The inputs
are the tower strip bounds and the exponentially summable path
expansions at the two adjacent stages. The passage to the physical
kernels does not require a small-field assumption.

\subsection{Counting types before periodization}
Let $K_{ab}(z)$ be a scalar cover kernel, with finite type sets
$\mathcal A,\mathcal B$, satisfying, for every fixed pair of types,
\[
 \sum_{z\in\mathbb Z^4}|K_{ab}(z)|e^{\eta|z|}\le C,
 \qquad \eta\ge0.
\]
A physical entry from $(x,a)$ to $(y,b)$ is a sum of selected terms
$K_{ab}(z)\operatorname{Ad}(U_{\gamma_{x,z}})$ with
$y=x+\operatorname{wrap}(z)$. Adjoint transport is an isometry of the
real colour space, including along winding paths, so the operator
norm of each term is at most $|K_{ab}(z)|$. Absolute convergence
therefore permits the following rearrangements:
\begin{equation}\label{eq:guide-periodic-rows}
 \sup_{x,a}\sum_{y,b}\|F_{(x,a),(y,b)}\|\le|\mathcal B|C,
 \qquad
 \sup_{y,b}\sum_{x,a}\|F_{(x,a),(y,b)}\|\le|\mathcal A|C.
\end{equation}
For the first, fix $z,b$: exactly one $y$ can contribute. For the
second, fix $z,a$: exactly one $x=y-\operatorname{wrap}(z)$ can
contribute. Dropping the truncation only increases the absolute sum.
Thus different lifts of one periodic displacement are all paid.

The wide kernel $\Phi$ has fine doubled plaquette rows and
coarse doubled plaquette columns. In coarse coordinates a fine type
includes its residue in the block. There are $L^4$ residues. Its
fixed-type strip cost is $C/L^2$, so
\eqref{eq:guide-periodic-rows} gives, after absorbing only fixed
orientation counts,
\begin{equation}\label{eq:guide-allowance}
 \begin{array}{c|cc}
 &\text{row sum}&\text{column sum}\\ \hline
 \Phi&a/L^2&aL^2\\
 \Phi'&a&a\\
 N&a&a
 \end{array}
\end{equation}
with $a$ fixed before $L$. Here $\Phi'$ has predictor-bond rows and coarse doubled-plaquette
columns; $N$ has auxiliary coarse rows and the same columns. Their fixed-type strip costs are bounded
independently of $L$. The tower estimates provide precisely these
three strip inputs; no residue factor is present in the narrow two
kernels. In particular the column bound for $\Phi$ cannot be replaced
by its row bound.

All three kernels are cut at path length at most
$R_0=\lfloor B/100000\rfloor$ in fine units. The endpoint of the
wide path beginning at a fine site of residue $r$ and ending at
coarse displacement $z$ is the embedded coarse point at parent plus
$z$: this is the path from $r$ to $Lz$ on the cover, reduced modulo
the fine period. The parent changes by at most one coarse step on
every fine bond. For the narrow paths the ordinary endpoint bound
applies. Thus every nonzero entry has anchor separation at most
$B/50$ in coarse units. This is an incidence assertion for these
kernels, and does not bound all auxiliary reads of a larger operation.

\subsection{Absolute path expansions and their transposes}
At one stage write the path expansion of a plaquette field in the form
\[
 \ell_p(U)=\sum_{i\ge0}c_{a,i}\,
       \operatorname{Ad}(U_{\gamma_{p,i}})F(U_{p+s_{a,i}}),
 \qquad p=(x,a).
\]
The label $i$ includes the target plaquette orientation and the
transport path. Its length $\ell_{a,i}$ bounds the distance from $x$
to the target anchor. The stage input is
\begin{equation}\label{eq:guide-path-input}
 \sum_i|c_{a,i}|e^{\sigma\ell_{a,i}}\le P,
 \qquad \sigma>0,\quad P>0,
\end{equation}
uniformly in $a$, with clipping relaxation radius $r_0>0$ satisfying
$\sigma r_0\ge1$. Both stages have such bounds, with constants fixed
before the step side $L$. We retain the full label set when different
labels land at the same periodic plaquette.

Keep terms with $d\ell_{a,i}\le R_0$, where $d=1$ on the fine stage
and $d=L$ on the coarse stage. Let $\ell_p^{\rm cut}$ be their sum
and let
\[
 k_{pu}=\sum_{\substack{i:d\ell_{a,i}\le R_0\\p+s_{a,i}=u}}
                  |c_{a,i}|.
\]
Absolute convergence follows from \eqref{eq:guide-path-input}.
The curvature bound $\|F(U_u)\|\le r(U_u)$ gives
\begin{equation}\label{eq:guide-path-sums}
 \|\ell_p^{\rm cut}(U)\|\le\sum_u k_{pu}r(U_u),\quad
 \sum_u k_{pu}\le P,\quad \sum_p k_{pu}\le N_pP.
\end{equation}
For the column bound, fix an input orientation $a$ and a label $i$.
Translation by its offset is a permutation of the periodic sites,
so at most one anchor can land at the specified $u$. Sum first over
labels and then over the $N_p$ possible input orientations. This
argument also explains why merging repeated targets before taking
absolute values would not justify the bound.

A retained site satisfies
$r(U_u)\le t\exp\{d(p,u)/r_0\}$ for every target $u$. Since
$d(p,u)\le\ell_{a,i}$ and $1/r_0\le\sigma$,
\eqref{eq:guide-path-input} implies
\begin{equation}\label{eq:guide-kept}
 \|\ell_p^{\rm cut}(U)\|\le Pt.
\end{equation}
The direct component is $\sqrt{a_0}F(U_p)$, where $a_0>0$ is the
stage quadratic coefficient. Its clipped version multiplies it by
$\min\{1,t/r(U_p)\}$, with multiplier $1$ at zero. Its norm is at
most $\sqrt{a_0}t$. The second component is
$\ell_p^{\rm cut}$ at retained sites and zero at discarded sites.
Thus the two-component clipped field $v$ satisfies the pointwise
bound
\begin{equation}\label{eq:guide-clipped}
 \|v_{p,\varepsilon}\|\le A_0t,
 \qquad A_0=(N_p+1)P+\sqrt{a_0}.
\end{equation}
This same $A_0$ bounds both row and column sums in
\eqref{eq:guide-path-sums}.

The selected and discarded rows have the locality needed for the
energy proof. For a cut path, the parent separation is at most its
length plus one, hence at most $R_0+1$. A witness to discarding at
lattice distance $k$ obeys $t e^{k/r_0}<r(U_u)\le2$. For
$t=P(h)/\sqrt h\ge h^{-1/2}$ this implies
$k<r_0\log(2\sqrt h)$. With $h=x+c_L$, $c_L$ fixed, this is at most
$r_0\log x$ eventually. Since $B=\lceil(\log x)^2\rceil$, both
$R_0+1$ and $k+1$ are at most $B/50$ eventually. The parent map
increases distances by at most one when compared with fine lattice
distance; therefore these bounds give the asserted coarse anchor
locality too. All tails are uniform in the period.

Finally let $s=t(x)$ be the coarse threshold and $t=t(x+c_L)$ the
fine threshold. The shift $c_L$ is nonnegative: its group factor is
a positive multiple of a sum of squares of adjoint matrix entries,
and $\log L\ge0$. For large $x$, $s\le2t$: the numerator
$P(x)=\exp\{(\log\log x)^2\}$ is increasing once $\log x\ge1$,
and $\sqrt{x+c_L}\le2\sqrt x$ for $x\ge c_L$.
Choose $A$ larger than the two stage constants $A_0$. Then
\eqref{eq:guide-clipped}, applied on the coarse stage, and
\eqref{eq:guide-allowance} imply
\begin{equation}\label{eq:guide-maximum}
 \|v_j\|\le dt,\qquad
 \|(\Phi v)_i\|\le dL^{-2}t,\qquad
 \|(\Phi'v)_b\|\le dt,
 \qquad d=2A(1+a).
\end{equation}
Indeed a matrix with operator-norm row sum $r$ sends a field bounded
by $u$ to a field bounded by $ru$. These are exactly the guide bounds
used for the fine and predictor energy rows in
Section~\ref{sec:seed-energy}, with constants fixed before $L,W$.
The derivation uses the stated tower and path-expansion inputs; the
construction of the path expansions remains a separate input to be
proved below.
